\section{Miary związane z całkowaniem}
Określenie „miara Lebesgue'a--Riemanna” nie jest standardową nazwą jednej konstrukcji. Rozróżniamy zawartość Jordana, miarę Lebesgue'a i miary Lebesgue'a--Stieltjesa.
\subsection*{Miara Jordana („riemannowska”)}
Dla ograniczonego $A\subset\R$ kładziemy $c_*(A)=\sup\{|E|:E\subset A,\ E\ \text{elementarny}\}$ oraz $c^*(A)=\inf\{|E|:A\subset E,\ E\ \text{elementarny}\}$, gdzie zbiory elementarne są skończonymi sumami zbiorów ograniczonych przedziałów, a $|E|$ to suma ich długości po sprowadzeniu do przedziałów o rozłącznych wnętrzach. Zbiór jest mierzalny w sensie Jordana, jeżeli $c_*(A)=c^*(A)$.
\begin{stwierdzenie}
Dla $A\subset[a,b]$ następujące warunki są równoważne: (i) $A$ jest mierzalny w sensie Jordana; (ii) $\ind A\in\mathcal R[a,b]$; (iii) $\lambda(\partial A)=0$. Wówczas $c(A)=\int_a^b\ind A(x)\,\mathrm dx=\lambda(A)$.
\end{stwierdzenie}
\begin{proof}
Niech $P$ będzie podziałem $[a,b]$. Suma domkniętych komórek zawartych w $A$ jest zbiorem elementarnym zawartym w $A$, a jej długość jest sumą dolną $L(\chi_A,P)$. Suma domkniętych komórek przecinających $A$ zawiera $A$, a jej długość jest sumą górną $U(\chi_A,P)$. Zatem
\[
 L(\chi_A,P)\le c_*(A)\le c^*(A)\le U(\chi_A,P).
\]
Jeżeli $\chi_A$ jest całkowalna Riemanna, różnica skrajnych wyrazów może być dowolnie mała, więc $c_*(A)=c^*(A)$, a wspólna wartość jest całką Riemanna.

Odwrotnie, jeżeli $c_*(A)=c^*(A)$, dla $\varepsilon>0$ wybieramy zbiory elementarne $E\subset A\subset H$ z $|H|-|E|<\varepsilon$. Każdy punkt brzegu $A$ leży w $H\setminus E$ albo w brzegu $E$ lub $H$: poza tymi zbiorami jest bowiem albo we wnętrzu $E$, albo na zewnątrz domknięcia $H$. Brzegi zbiorów elementarnych są skończone, a $H\setminus E$ jest elementarny i ma długość $|H|-|E|$. Stąd $\lambda_1(\partial A)\le |H|-|E|<\varepsilon$. Dowolność $\varepsilon$ daje (iii).
Punkty nieciągłości $\chi_A$ tworzą względny brzeg $A$ w $[a,b]$, różniący się od brzegu w $\R$ co najwyżej końcami $a,b$. Z kryterium Lebesgue'a całkowalności Riemanna wynika (iii) $\Rightarrow$ (ii).
Jeżeli brzeg jest zerowy, zbiór $A$ różni się od swojego wnętrza o podzbiór brzegu, więc jest mierzalny w sensie Lebesgue'a. Dla dowolnych $E\subset A\subset H$ jak wyżej zachodzi $|E|\le\lambda_1(A)\le |H|$. Przechodząc do kresów, otrzymujemy $c(A)=\lambda_1(A)$.
\end{proof}
Zbiory mierzalne w sensie Jordana tworzą algebrę, nie $\sigma$-algebrę: $\Q\cap[0,1]$ jest przeliczalną sumą zbiorów jednopunktowych mierzalnych w sensie Jordana, ale sam nie jest mierzalny w tym sensie. Przejście od skończonych do przeliczalnych pokryć prowadzi do miary zewnętrznej Lebesgue'a $\lambda^*(A)=\inf\{\sum_k|I_k|:A\subset\bigcup_kI_k\}$ oraz konstrukcji Carathéodory'ego. Dla ograniczonego $A$ mamy $c_*(A)\le\lambda_*(A)\le\lambda^*(A)\le c^*(A)$.

\subsection*{Miary Lebesgue'a--Stieltjesa}
\begin{twierdzenie}[Konstrukcja miary Lebesgue'a--Stieltjesa]\label{it:thm:stieltjes-construction}
Niech $\alpha:\R\to\R$ będzie niemalejąca i prawostronnie ciągła. Istnieje dokładnie jedna lokalnie skończona miara borelowska $\mu_\alpha$ na $\R$, dla której
\[
 \mu_\alpha((u,v])=\alpha(v)-\alpha(u)\qquad(u<v).
\]
Jej uzupełnienie nazywamy miarą Lebesgue'a--Stieltjesa. Odwrotnie, każda lokalnie skończona miara borelowska $\nu$ na $\R$ ma tę postać. Funkcja $\alpha$ jest wyznaczona z dokładnością do stałej.
\end{twierdzenie}
\begin{proof}
Na pierścieniu $\mathcal A$ skończonych sum rozłącznych ograniczonych przedziałów $(u,v]$ definiujemy $m$ jako sumę przyrostów $\alpha$. Wspólne rozdrobnienie skończonych rodzin w ich końcach oraz teleskopowanie przyrostów dowodzą niezależności od przedstawienia i skończonej addytywności. W szczególności $m$ jest nieujemna, monotoniczna i skończenie podaddytywna.

Udowodnimy ciągłość z góry w zbiorze pustym. Niech $A_n\in\mathcal A$, $A_n\downarrow\varnothing$ i niech $\varepsilon>0$. Każdy przedział $(u,v]$ w skończonym przedstawieniu $A_n$ zastępujemy zwartym podprzedziałem $[u+h,v]$, $0<h<v-u$. Na mocy prawostronnej ciągłości $\alpha$ wybieramy $h$ tak, by suma wartości $m((u,u+h])$ dla wszystkich komórek $A_n$ była mniejsza od $\varepsilon 2^{-n}$. Otrzymujemy zbiór zwarty $K_n\subset A_n$ oraz $D_n\in\mathcal A$ z $A_n\setminus K_n\subset D_n$ i $m(D_n)<\varepsilon 2^{-n}$. Dla pustego $A_n$ bierzemy $K_n=D_n=\varnothing$.
Ponieważ $\bigcap_n K_n\subset\bigcap_n A_n=\varnothing$, własność skończonych przecięć zbiorów zwartych daje $N$ takie, że $\bigcap_{n=1}^N K_n=\varnothing$. Dla $x\in A_N$ mamy $x\in A_n$ dla każdego $n\le N$, a $x\notin K_n$ dla pewnego takiego $n$. Zatem $A_N\subset\bigcup_{n=1}^N D_n$ i
\[
 m(A_N)\le \sum_{n=1}^Nm(D_n)<\varepsilon.
\]
Monotoniczność daje $m(A_n)\to0$.

Jeśli $A=\bigsqcup_{n\ge1}B_n\in\mathcal A$ oraz $B_n\in\mathcal A$, reszty $R_N=A\setminus\bigcup_{n=1}^N B_n$ należą do $\mathcal A$ i maleją do zbioru pustego. Skończona addytywność daje $m(A)=\sum_{n=1}^N m(B_n)+m(R_N)$; przejście do granicy dowodzi, że $m$ jest przedmiarą. Twierdzenie Carathéodory'ego o rozszerzaniu przedmiary daje miarę na $\sigma(\mathcal A)=\mathcal B(\R)$. Jednoznaczność wynika z $\sigma$-skończoności: przedziały $(-n,n]$ pokrywają prostą i mają skończone masy. Każdy zbiór zwarty leży w jednym z tych przedziałów, więc miara jest lokalnie skończona.

Odwrotnie, dla lokalnie skończonej $\nu$ definiujemy
\[
 \alpha(x)=
 \begin{cases}
 \nu((0,x]),&x\ge0,\\
 -\nu((x,0]),&x<0.
 \end{cases}
\]
W obu przypadkach $\alpha(v)-\alpha(u)=\nu((u,v])$, także dla $u<0\le v$ przez addytywność. Z nieujemności miary wynika monotoniczność. Dla $y\downarrow x$ mamy $\nu((x,y])\to0$ z ciągłości miary z góry na ograniczonym przedziale, a zatem $\alpha$ jest prawostronnie ciągła. Jednoznaczność miary daje $\nu=\mu_\alpha$. Jeżeli dwie funkcje mają te same przyrosty, ich różnica jest stała.
\end{proof}
Dla $\alpha(x)=x$ otrzymujemy $\lambda_1$, a dla $\alpha=\chi_{[c,\infty)}$ miarę Diraca $\delta_c$, ponieważ obie spełniają właściwy wzór na przedziałach $(u,v]$. Dla integratora prawostronnie ciągłego w porównaniu z całką Riemanna--Stieltjesa na $[a,b]$ ograniczamy miarę do $(a,b]$ (zob. twierdzenie \ref{it:thm:rs-oscillation} i stwierdzenie \ref{st:ks-miara-konwencja}).
\subsection*{Rozkład miary i całkowanie względem gęstości}
% Import TM-02, TM-17, TM-18; zachowane etykiety źródłowe.
\begin{definicja}[Bezwzględna ciągłość miary]
\label{def:ac-measure-17}
Niech $\mu$ i $\nu$ będą miarami nieujemnymi na $(X,\mathcal F)$.
Mówimy, że $\nu$ jest \emph{bezwzględnie ciągła względem $\mu$},
i piszemy $\nu\ll\mu$, jeżeli dla każdego $E\in\mathcal F$
z warunku $\mu(E)=0$ wynika $\nu(E)=0$.
Dla skończonej miary zespolonej $\eta$ przyjmujemy tę samą definicję,
równoważnie $|\eta|\ll\mu$, gdzie $|\eta|$ jest wariacją całkowitą.
\end{definicja}
Miary nieujemne $\nu,\mu$ na wspólnej przestrzeni mierzalnej są osobliwe ($\nu\perp\mu$), gdy istnieje mierzalny zbiór $S$ taki, że $\mu(S)=0$ i $\nu(X\setminus S)=0$. Korzystamy z twierdzenia o rozkładzie Lebesgue'a oraz twierdzenia Radona--Nikodyma (Teoria miary, rozdziały 17--18).
\begin{definicja}[Atom i miara bezatomowa]
\label{def:atom}
Niech $(X,\mathcal F,\mu)$ będzie przestrzenią z miarą.
Zbiór $A\in\mathcal F$ nazywamy \emph{atomem miary $\mu$},
jeżeli $\mu(A)>0$ oraz dla każdego $B\in\mathcal F$
zawartego w $A$ zachodzi $\mu(B)=0$ lub
$\mu(A\setminus B)=0$.
Miarę $\mu$ nazywamy \emph{bezatomową}, jeśli nie ma atomów.
\end{definicja}
\begin{twierdzenie}[Dalszy rozkład miary osobliwej na prostej]
\label{thm:atomic-cont-singular-17}
Niech $\nu$ będzie skończoną miarą borelowską na
$\mathbb R$. Zbiór
$A=\{x\in\mathbb R:\nu(\{x\})>0\}$ jest co najwyżej
przeliczalny, a miary
\[
 \nu_{\mathrm{at}}(E):=\nu(E\cap A),\qquad
 \nu_{\mathrm{na}}(E):=\nu(E\setminus A)
 \quad(E\in\mathcal B(\mathbb R))
\]
są jednoznacznie określone. Pierwsza ma postać
$\nu_{\mathrm{at}}=\sum_{a\in A}\nu(\{a\})\delta_a$,
a druga jest bezatomowa. Jeśli $\nu\perp\lambda_1$,
to również $\nu_{\mathrm{at}}\perp\lambda_1$ i
$\nu_{\mathrm{na}}\perp\lambda_1$.
\end{twierdzenie}
\begin{proof}
Dla każdego $n\ge1$ zbiór
$A_n=\{x:\nu(\{x\})\ge1/n\}$ jest skończony:
gdyby zawierał $m$ różnych punktów, mielibyśmy
$\nu(\mathbb R)\ge m/n$. Ponieważ
$A=\bigcup_{n\ge1}A_n$, jest przeliczalny.
Przeliczalna addytywność daje podany wzór na
$\nu_{\mathrm{at}}$ oraz rozkład $\nu$.
Z definicji $\nu_{\mathrm{na}}(\{x\})=0$ dla
każdego $x$. Wykażemy, że miara ta nie ma także
żadnych atomów w sensie teorii miary.

Przypuśćmy, że mierzalny zbiór $B$ jest atomem
$\nu_{\mathrm{na}}$ o dodatniej masie $b$:
każdy mierzalny podzbiór $B$ ma masę $0$ lub $b$.
Ponieważ $\mathbb R=\bigcup_{m\ge1}[-m,m]$,
pewien domknięty ograniczony przedział $I_0$
spełnia $\nu_{\mathrm{na}}(B\cap I_0)=b$.
Dzielimy $I_0$ na dwa domknięte przedziały równej
długości. Co najmniej jeden z nich, nazwany $I_1$,
ma $\nu_{\mathrm{na}}(B\cap I_1)>0$, a więc masę $b$.
Powtarzając, otrzymujemy zstępujący ciąg przedziałów
zamkniętych $I_n$ o długościach malejących do zera,
dla których $\nu_{\mathrm{na}}(B\cap I_n)=b$.
Ich przecięcie jest singletonem $\{x\}$.
Z ciągłości skończonej miary z góry mamy
$\nu_{\mathrm{na}}(B\cap\{x\})=b>0$,
sprzecznie z zerowością wszystkich singletonów.
Zatem $\nu_{\mathrm{na}}$ jest bezatomowa.

Jeśli $\nu$ skupia się na zbiorze borelowskim $S$
o $\lambda_1(S)=0$, to obie części jako ograniczenia
$\nu$ również skupiają się na $S$, więc są osobliwe.
\end{proof}
\begin{lemat}[Całkowanie względem miary o gęstości]
\label{lem:integration-density-18}
Niech $\lambda$ będzie miarą nieujemną na $(X,\mathcal F)$,
$w:X\to[0,\infty]$ funkcją mierzalną oraz
$\mu(E)=\int_Ew\,\mathrm d\lambda$ dla $E\in\mathcal F$.
Wówczas dla każdej mierzalnej $\varphi:X\to[0,\infty]$
zachodzi
\begin{equation}\label{eq:integration-density-18}
\int_X\varphi(x)\,\mathrm d\mu(x)
=\int_X\varphi(x)w(x)\,\mathrm d\lambda(x),
\end{equation}
przy konwencji $0\cdot\infty=0$ dla iloczynu
nieujemnych funkcji mierzalnych. Jeśli $\varphi$ jest
zespolona i całkowalna względem $\mu$, to analogiczna
równość zachodzi dla całek zespolonych.
\end{lemat}
\begin{proof}
Dla $\varphi=\chi_A$ jest to definicja $\mu$.
Z liniowości dla funkcji prostych nieujemnych otrzymujemy
równość dla wszystkich takich funkcji. Każdą funkcję
$\varphi\ge0$ aproksymujemy z dołu ciągiem mierzalnych
funkcji prostych $\varphi_n\uparrow\varphi$.
Dwukrotne zastosowanie twierdzenia Beppo Leviego
daje \eqref{eq:integration-density-18}. W przypadku
funkcji zespolonej stosujemy tę tożsamość do części
dodatnich i ujemnych jej części rzeczywistej i urojonej,
które są całkowalne dzięki założeniu.
\end{proof}

\begin{twierdzenie}[Kryterium oparte na rozkładzie miary]
Niech $\mu_\alpha$ będzie skończoną miarą Lebesgue'a--Stieltjesa ograniczoną do $(a,b]$. Istnieje rozkład
\[
\mu_\alpha=w\lambda_1+\mu_{\mathrm{sc}}+\sum_{k\in J}c_k\delta_{t_k},
\]
gdzie wszystkie miary rozpatrujemy na $(a,b]$, $w\ge0$ jest mierzalna, $\mu_{\mathrm{sc}}\perp\lambda_1$ jest bezatomowa, $J$ jest zbiorem co najwyżej przeliczalnym, $t_k$ są różne, $c_k>0$ i $\sum_{k\in J}c_k<\infty$. Dla każdej funkcji borelowskiej $f:(a,b]\to\R$ zachodzi
\[
\int_{(a,b]}|f|\,\mathrm d\mu_\alpha
=\int_{(a,b]}|f|w\,\mathrm d\lambda_1+
\int_{(a,b]}|f|\,\mathrm d\mu_{\mathrm{sc}}+
\sum_{k\in J}c_k|f(t_k)|.
\]
Funkcja jest bezwzględnie całkowalna względem $\mu_\alpha$ wtedy i tylko wtedy, gdy prawa strona jest skończona.
\end{twierdzenie}
\begin{proof}
Obie miary $\mu_\alpha$ i $\lambda_1|_{(a,b]}$ są skończone, więc twierdzenie o rozkładzie Lebesgue'a daje część bezwzględnie ciągłą i osobliwą. Twierdzenie Radona--Nikodyma daje nieujemną gęstość $w$ pierwszej części względem $\lambda_1$. Drugą część przedłużamy przez zero na $\R$ i stosujemy twierdzenie \ref{thm:atomic-cont-singular-17}. Masa atomów sumuje się do masy części atomowej, nie większej od $\mu_\alpha((a,b])$.
Dla nieujemnej borelowskiej funkcji prostej równość całek względem sumy miar wynika z definicji całki i przeliczalnej addytywności sum nieujemnych. Dla dowolnej nieujemnej funkcji mierzalnej przybliżamy ją od dołu rosnącym ciągiem funkcji prostych i stosujemy twierdzenie Beppo Leviego do każdej miary oraz do sumy po $J$ (miara licząca). Lemat \ref{lem:integration-density-18} identyfikuje całkę względem $w\lambda_1$, a całka względem $\delta_{t_k}$ jest wartością w $t_k$. Stosujemy otrzymaną równość do $|f|$.
\end{proof}

\subsection*{Miara Cantora: obliczenia i porównanie całek}
\begin{przyklad}[Miara Cantora — osobliwość bez atomów]
\label{ex:cantor-measure-17}
Podamy konstrukcję niewymagającą wcześniejszego twierdzenia
o miarach Lebesgue'a--Stieltjesa. Dla $x\in[0,1)$ niech
\[
b_n(x)=\lfloor2^nx\rfloor-2\lfloor2^{n-1}x\rfloor\in\{0,1\},
\qquad n\ge1.
\]
Są to cyfry rozwinięcia dwójkowego, przy czym dla liczb
dwójkowo wymiernych przyjęliśmy rozwinięcie zakończone
samymi zerami. Określamy
$T(x)=\sum_{n\ge1}2b_n(x)3^{-n}$ dla $x\in[0,1)$
oraz $T(1)=1$. Każda funkcja $b_n$ jest borelowska,
a szereg definiujący $T$ jest jednostajnie zbieżny,
więc $T:[0,1]\to[0,1]$ jest borelowska. Jej wartości
należą do środkowego zbioru Cantora $C$, ponieważ ich
rozwinięcia trójkowe używają wyłącznie cyfr $0$ i $2$.

Niech $\kappa=T_*(\lambda_1|_{[0,1]})$. Jest to miara
probabilistyczna, $\kappa(C)=1$. Na poziomie $n$ konstrukcji zbioru Cantora
każdy z $2^n$ rozłącznych przedziałów podstawowych
$I_{\varepsilon_1,\ldots,\varepsilon_n}$ odpowiada
słowu $(\varepsilon_1,\ldots,\varepsilon_n)\in\{0,1\}^n$:
jego punkty mają rozwinięcie trójkowe rozpoczynające się
od cyfr $2\varepsilon_1,\ldots,2\varepsilon_n$.
Z wyjątkiem przeliczalnego zbioru końców przedziałów dwójkowych
przeciwobraz $I_{\varepsilon_1,\ldots,\varepsilon_n}$
przez $T$ jest zbiorem liczb, których pierwsze $n$
cyfr dwójkowych to $\varepsilon_1,\ldots,\varepsilon_n$.
Jest to przedział dwójkowy długości $2^{-n}$, więc
$\kappa(I_{\varepsilon_1,\ldots,\varepsilon_n})=2^{-n}$.
Każdy $y\in C$ należy na każdym poziomie do jednego
takiego przedziału, zatem
$\kappa(\{y\})\le2^{-n}$ dla wszystkich $n$ i
$\kappa(\{y\})=0$. Dla $y\notin C$ mamy
$\kappa(\{y\})=0$, gdyż $\kappa$ jest skupiona na $C$.
Miara $\kappa$ jest więc bezatomowa. Ponieważ
$\lambda_1(C)=0$, otrzymujemy $\kappa\perp\lambda_1$.
Pokazuje to, że \emph{osobliwość} nie jest tym samym co
\emph{dyskretność}. Dystrybuanta miary $\kappa$ jest
klasyczną funkcją Cantora.
\end{przyklad}
W dalszych przykładach oznaczamy funkcję Cantora przez $C$, środkowotrzeci zbiór Cantora przez $K$, a skonstruowaną miarę przez $\mu_C=\kappa$.
\begin{przyklad}[Momenty miary osobliwej]
Dla losowej zmiennej $X$ o rozkładzie $\mu_C$ mamy $X\stackrel d= (B+X')/3$, gdzie $B$ przyjmuje wartości $0,2$ z prawdopodobieństwami $1/2$, a $X'$ jest niezależną kopią $X$. Wynika to z tego, że dla miary Lebesgue'a każdy skończony ciąg cyfr dwójkowych długości $n$ ma prawdopodobieństwo $2^{-n}$. Cyfry są więc niezależne i każda ma rozkład Bernoulliego z parametrem $1/2$; po oddzieleniu pierwszej cyfry szereg pozostałych cyfr ma po przeskalowaniu ten sam rozkład i jest niezależny od pierwszej cyfry. Stąd
$\mathbb E X=\frac13\mathbb E X+\frac13$, czyli $\mathbb E X=1/2$, a także $\mathbb E X^2=\frac19\mathbb E X^2+\frac29\mathbb E X+\frac29$, czyli $\mathbb E X^2=3/8$. Zatem
\[
\int_0^1 1\,\mathrm dC=1,\qquad \int_0^1 x\,\mathrm dC=\frac12,\qquad \int_0^1 x^2\,\mathrm dC=\frac38.
\]
Ponadto z całkowania przez części, ponieważ $C$ jest ciągła i ma wahanie ograniczone, wynika $2\int_0^1 C\,\mathrm dC=C(1)^2-C(0)^2=1$. Choć $C'=0$ prawie wszędzie, nie wolno zastępować $\mathrm dC$ przez $C'(x)\,\mathrm dx$: zgubilibyśmy część osobliwą miary. Ponieważ $x\mapsto x^2$ jest ciągła, równości obowiązują zarówno dla całek Riemanna--Stieltjesa, jak i Lebesgue'a--Stieltjesa.
\end{przyklad}
\begin{przyklad}[Ta sama funkcja, różne miary]
Niech $K$ będzie zbiorem Cantora. Funkcja $\ind K$ jest całkowalna Riemanna względem $\lambda$ z całką $0$, ponieważ $K$ ma miarę Lebesgue'a zero. Jest również całkowalna Lebesgue'a--Stieltjesa względem $\mu_C$ z całką $1$. Natomiast zwykła całka Riemanna--Stieltjesa $\int_0^1\ind K\,\mathrm dC$ nie istnieje: punkty nieciągłości funkcji $\ind K$ tworzą $K$, a $\mu_C(K)=1$, więc nie jest spełnione kryterium oscylacyjne.
\end{przyklad}

% ============================================================
% UZUPELNIENIE MATEMATYCZNE: ROZDZIAL 10
% ============================================================
\subsection*{Atomy i wartości na końcach: porównanie z KS}
\begin{stwierdzenie}[Reprezentacja całki KS przez miarę przy prawostronnej ciągłości]\label{st:ks-miara-konwencja}
Niech $\alpha:[a,b]\to\R$ będzie niemalejąca, prawostronnie ciągła na $[a,b)$ i niech $\mu_\alpha$ będzie miarą na $(a,b]$ zdefiniowaną przez $\mu_\alpha((s,t])=\alpha(t)-\alpha(s)$. Jeśli $f:[a,b]\to\R$ jest regulowana, to całka Kurzweila--Stieltjesa $\int_a^b f\,\mathrm d\alpha$ istnieje i
\[
\int_a^b f\,\mathrm d\alpha=\int_{(a,b]}f\,\mathrm d\mu_\alpha.
\]
W szczególności $\mu_\alpha(\{c\})=\alpha(c)-\alpha(c-)$ dla $a<c\le b$, a punkt $a$ nie wnosi atomu w tej konwencji.
\end{stwierdzenie}
\begin{proof}
Z twierdzenia \ref{tw:ks-measure} wiemy, że dla integrandu regulowanego i integratora BV całka istnieje oraz że zgadza się z całką Lebesgue'a--Stieltjesa w przyjętej konwencji. Tutaj $\alpha$ jest monotoniczna, więc ma wahanie ograniczone. Przechodząc z $s\uparrow c$ w równości $\mu_\alpha((s,c])=\alpha(c)-\alpha(s)$, otrzymujemy wskazaną masę atomu.
\end{proof}
\begin{przyklad}[Skok prawostronny i wartość w punkcie]
Dla $a<c<b$ zdefiniujmy $g(x)=\ind{(c,b]}(x)$. Funkcja $g$ nie jest prawostronnie ciągła w $c$, a jej przyrost na przedziale $[c,t]$, $t>c$, wynosi $1$. Dla dowolnej regulowanej funkcji $f$ całka KS istnieje i $\int_a^b f\,\mathrm dg=f(c)$ (jest to wkład prawego skoku). Zastąpienie $g$ przez prawostronnie ciągłą funkcję $g^+(x)=\ind{[c,b]}(x)$ nie zmienia tu wartości całki KS dla regulowanej $f$, lecz przenosi skok z prawej strony $c$ na lewą. Zwykła miara Stieltjesa jest związana z wybranym prawostronnie ciągłym reprezentantem.
\end{przyklad}
\begin{przyklad}[Dwa skoki przeciwnego znaku w jednym punkcie]
Niech $g(x)=\ind{\{c\}}(x)$, gdzie $a<c<b$. Funkcja ma skok lewy $\Delta^-g(c)=1$ i prawy $\Delta^+g(c)=-1$. Dla regulowanej $f$ wzór skokowy z rozdziału 6 daje $\int_a^b f\,\mathrm dg=f(c)-f(c)=0$. Nie należy utożsamiać takiej funkcji $g$ z dystrybuantą dodatniej miary: $g$ nie jest monotoniczna ani prawostronnie ciągła.
\end{przyklad}
\begin{przyklad}[Dwa różne końce przedziału]
Na $[a,b]$ niech $g_a=\ind{(a,b]}$ oraz $g_b=\ind{\{b\}}$. Dla dowolnej regulowanej $f$ zachodzi
$\int_a^b f\,\mathrm dg_a=f(a)$ oraz $\int_a^b f\,\mathrm dg_b=f(b)$ w sensie KS. Są to odpowiednio skok prawostronny w $a$ i skok lewostronny w $b$. W miarze $\mu_\alpha$ na $(a,b]$ pierwszy z nich nie występuje, ponieważ funkcja $g_a$ nie jest prawostronnie ciągła w $a$; drugi odpowiada atomowi w $b$.
\end{przyklad}
